\documentclass[12pt,a4paper]{article}

\usepackage[utf8]{inputenc}
\usepackage[T1]{fontenc}
\usepackage{lmodern}
\usepackage[spanish,es-tabla]{babel}
\usepackage[margin=2.5cm]{geometry}
\usepackage{setspace}
\usepackage{booktabs}
\usepackage{graphicx}
\usepackage{enumitem}
\usepackage{titlesec}
\usepackage{parskip}
\usepackage{xcolor}
\usepackage{hyperref}
\usepackage{array}
\usepackage{tabularx}

\setstretch{1.15}
\setlength{\parskip}{0.6em}
\setlength{\parindent}{0pt}

\definecolor{graytitle}{gray}{0.20}
\definecolor{graysoft}{gray}{0.45}

\titleformat{\section}
  {\normalfont\large\bfseries\color{graytitle}}
  {\thesection.}{0.6em}{}
\titlespacing*{\section}{0pt}{1.6em}{0.8em}

\titleformat{\subsection}
  {\normalfont\normalsize\bfseries\color{graytitle}}
  {\thesubsection}{0.5em}{}
\titlespacing*{\subsection}{0pt}{1.2em}{0.4em}

\hypersetup{
  colorlinks=true,
  linkcolor=black,
  urlcolor=graytitle,
  citecolor=black,
  pdftitle={Herramientas internas y plataforma de infraestructura en CrediFlexi},
  pdfauthor={Uzziel Valdez Ordonez}
}

\setlist[itemize]{leftmargin=1.4em, itemsep=0.2em, topsep=0.3em}
\setlist[enumerate]{leftmargin=1.6em, itemsep=0.2em, topsep=0.3em}

\renewcommand{\arraystretch}{1.25}

\begin{document}

% ============================================================
% PORTADA
% ============================================================
\begin{titlepage}
\centering
\vspace*{3.5cm}

{\Large\bfseries Herramientas internas y plataforma\par}
\vspace{0.3em}
{\Large\bfseries de infraestructura en CrediFlexi\par}

\vspace{1.4cm}

{\large\color{graysoft} Propuesta técnica y plan de adopción\par}

\vspace{2.8cm}

\rule{6cm}{0.3pt}

\vspace{1em}

{\normalsize Mesa de Tickets\par}
\vspace{0.3em}
{\normalsize Calculadora de Score Crediticio\par}

\vspace{1em}

\rule{6cm}{0.3pt}

\vspace{2.8cm}

{\normalsize Uzziel Valdez Ordoñez\par}
\vspace{0.3em}
{\small\color{graysoft} Becario de Data Science e Ingeniería de Software\par}

\vfill

{\small 4 de mayo de 2026 \quad $\cdot$ \quad Versión 1.0\par}

\end{titlepage}

% ============================================================
% TOC
% ============================================================
\thispagestyle{empty}
\tableofcontents
\newpage

% ============================================================
% 1. RESUMEN EJECUTIVO
% ============================================================
\section{Resumen ejecutivo}

Este documento presenta dos herramientas internas ya construidas y operativas en versión prototipo:
una \textbf{Mesa de Tickets} para la gestión de incidencias y solicitudes internas,
y una \textbf{Calculadora de Score Crediticio} para la evaluación de clientes.

Más allá del valor individual de cada herramienta,
el objetivo central de esta propuesta es \textbf{formalizar el stack tecnológico que las soporta}
como plataforma común de desarrollo interno de CrediFlexi.
Una vez definida y bajo gobernanza del Área de Sistemas,
esta plataforma permite construir herramientas adicionales,
internas o comerciales,
con bajo costo marginal y sin reabrir cada vez la conversación de infraestructura.

El stack propuesto es estándar de la industria, sin tecnologías propietarias,
y portable a otras plataformas si en el futuro fuese necesario.
El costo de operación es \textbf{cero durante el piloto}
y se ubica entre \textbf{\$360 y \$810 MXN mensuales} en producción,
para todas las herramientas que se construyan sobre esta base.

La decisión solicitada es la \textbf{aprobación para iniciar la migración
de las cuentas técnicas, hoy a título personal, hacia infraestructura corporativa}
bajo administración del Área de Sistemas.

% ============================================================
% 2. PROBLEMA ACTUAL
% ============================================================
\section{Problema actual}

Las herramientas internas operativas viven hoy en \textbf{Google Apps Script}
sobre hojas de cálculo de Google.
Esta arquitectura cumplió su propósito inicial
pero presenta limitaciones crecientes conforme aumenta el volumen y la complejidad operativa:

\begin{itemize}
  \item Control de acceso básico, sin segmentación clara por área o rol.
  \item Trazabilidad limitada: no es posible reconstruir con certeza quién hizo qué y cuándo.
  \item Degradación de rendimiento con volumen creciente de registros.
  \item Mantenimiento frágil: cambios sencillos requieren intervención manual sobre las hojas.
  \item Conocimiento concentrado en personas: la rotación implica pérdida operativa.
  \item No existe una vía sencilla para exponer estas herramientas a clientes
        ni integrarlas con sistemas corporativos.
\end{itemize}

Casos sentidos hoy:

\begin{itemize}
  \item \textbf{Mesa de tickets}: las incidencias internas se gestionan entre WhatsApp,
        correos y una hoja compartida. No existen métricas, SLA, ni auditoría formal.
  \item \textbf{Cálculo de score crediticio}: el flujo depende de un libro de Excel y scripts
        en Apps Script con parámetros dispersos, sin historial estructurado de evaluaciones.
\end{itemize}

% ============================================================
% 3. SOLUCIÓN PROPUESTA
% ============================================================
\section{Solución propuesta}

\subsection*{Las dos herramientas}

\textbf{Mesa de Tickets.}
Aplicación web para la gestión interna de incidencias, solicitudes
y soporte operativo entre áreas.
Centraliza en un único lugar lo que hoy vive disperso entre WhatsApp,
correos y hojas de cálculo.
Cualquier colaborador puede levantar un ticket,
asignarlo al área correspondiente,
darle seguimiento por hilos de respuesta, adjuntar evidencias
y cerrar el caso con confirmación dual:
del responsable técnico y del solicitante.
Cada acción queda registrada con autor, fecha y contexto,
lo que permite construir métricas operativas reales
---tiempo de respuesta, carga por área, casos vencidos---
y conservar el conocimiento institucional incluso ante rotación de personal.

Disponible como prototipo funcional en
\href{https://mesa-tickets.vercel.app}{\texttt{mesa-tickets.vercel.app}}.

\textbf{Calculadora de Score Crediticio.}
Aplicación web para la evaluación del puntaje crediticio de clientes.
Reemplaza el flujo actual basado en Excel y Apps Script
por un sistema con reglas y parámetros centralizados,
versionados y modificables sin tocar el código de la aplicación.
Cada evaluación queda guardada con su contexto completo:
qué evaluador la ejecutó, qué parámetros estaban vigentes en ese momento,
qué resultado se obtuvo y sobre qué cliente.
Esto entrega criterio uniforme entre evaluadores,
historial estructurado para análisis posteriores
y trazabilidad ante auditorías.

Ambas herramientas comparten la misma base técnica,
descrita en la sección de arquitectura.

\subsection*{El valor estructural: una plataforma reutilizable}

El diferenciador real de esta propuesta no son las dos herramientas en sí mismas,
sino el hecho de que ambas validan en producción un mismo stack técnico.
Una vez que ese stack esté formalmente bajo gobernanza corporativa,
construir nuevas herramientas implica solo desarrollo,
no decisiones de infraestructura.

Esto habilita a CrediFlexi a desplegar
\textbf{herramientas sin un límite definido},
tanto para uso \textbf{interno} ---socios, colaboradores y áreas operativas---
como para uso \textbf{comercial} ---productos, portales y servicios
de cara a clientes finales---.
El costo y el tiempo de cada nueva herramienta se reducen sustancialmente
porque la infraestructura, la autenticación, la gobernanza
y las prácticas de operación ya están resueltas.

\begin{figure}[h]
\centering
\includegraphics[width=0.92\textwidth]{plataforma.png}
\caption*{\small\color{graysoft} Figura 1. La plataforma común como base
de las herramientas actuales y futuras.}
\end{figure}

Ejemplos de iniciativas que podrían construirse sobre esta base:

\textbf{Para uso interno (socios y colaboradores):}
dashboards operativos por área, automatizaciones de procesos hoy manuales,
formularios estructurados con auditoría, integraciones con sistemas existentes vía API.

\textbf{Para uso comercial (cara a clientes):}
portales de autoservicio, calculadoras y simuladores públicos de productos crediticios,
formularios de solicitud o precalificación en línea, micrositios para campañas.

% ============================================================
% 4. ALCANCE DEL MVP
% ============================================================
\section{Alcance del MVP}

\begin{tabularx}{\textwidth}{@{}>{\raggedright\arraybackslash}X >{\raggedright\arraybackslash}X@{}}
\toprule
\textbf{Incluido en esta entrega} & \textbf{Fuera de alcance (fases posteriores)} \\
\midrule
Mesa de tickets funcional & Aplicación móvil nativa \\
Calculadora de score funcional & Integración con CRM corporativo \\
Inicio de sesión con Google Workspace & Portal externo orientado a clientes \\
Roles y permisos por área & Reportería avanzada (BI) \\
Auditoría y métricas operativas & Automatizaciones de procesos complejos \\
Hosting y subdominio corporativo & Operación multi-empresa \\
Documentación técnica y de operación & Migración de herramientas existentes \\
\bottomrule
\end{tabularx}

% ============================================================
% 5. ARQUITECTURA TÉCNICA
% ============================================================
\section{Arquitectura técnica}

El sistema opera con cuatro componentes independientes,
cada uno cumple un rol claro y puede ser sustituido sin afectar a los demás.

\begin{figure}[h]
\centering
\includegraphics[width=0.9\textwidth]{arquitectura.png}
\caption*{\small\color{graysoft} Figura 2. Arquitectura del stack común a las dos herramientas.}
\end{figure}

\subsection*{Componentes}

\begin{itemize}
  \item \textbf{GitHub} --- repositorio del código fuente.
        Único origen de verdad. Cada cambio queda versionado, atribuido y revisable.
        Permite revisión por pares antes de cualquier despliegue.
  \item \textbf{Vercel} --- plataforma de hosting.
        Despliega automáticamente cualquier cambio aprobado en GitHub.
        Provee HTTPS, escalado automático, red de distribución global,
        métricas, registros de actividad y reversión inmediata ante incidentes.
  \item \textbf{Supabase} --- base de datos y servicios de aplicación.
        Por debajo opera \emph{PostgreSQL} estándar, una de las bases de datos relacionales
        más adoptadas del mundo.
        Provee almacenamiento de archivos, autenticación
        y reglas de seguridad a nivel de fila (Row Level Security).
  \item \textbf{Google Workspace} --- proveedor de identidad corporativa.
        El inicio de sesión se delega vía OAuth 2.0.
        El control de quién puede acceder se mantiene en el panel de Workspace de la empresa.
        Al desactivar la cuenta de un colaborador, el acceso a las herramientas se revoca
        de forma automática.
\end{itemize}

\subsection*{Portabilidad y ausencia de \emph{vendor lock-in}}

Las decisiones técnicas se tomaron priorizando portabilidad.
PostgreSQL es un estándar abierto y exportable.
La aplicación corre en cualquier infraestructura compatible con Node.js.
El código vive en un repositorio independiente del proveedor de hosting.
Una eventual migración hacia infraestructura propia,
o hacia otro proveedor (AWS, Azure, servidores en oficinas),
es viable sin reescritura, en plazos de semanas y no de meses.

\subsection*{Seguridad}

\begin{itemize}
  \item Tráfico encriptado de extremo a extremo (HTTPS) en todas las comunicaciones.
  \item Autenticación delegada a Google Workspace; la aplicación no almacena contraseñas.
  \item Reglas de control de acceso aplicadas \emph{a nivel base de datos} (RLS),
        independientes del frontend.
  \item Vercel y Supabase cuentan con certificación SOC 2 Type II.
\end{itemize}

% ============================================================
% 6. COSTOS Y ESCENARIOS
% ============================================================
\section{Costos y escenarios}

Los costos siguientes consideran un tipo de cambio referencial de \$18 MXN por dólar
y aplican al stack completo, no por herramienta.

\begin{tabularx}{\textwidth}{@{}l l l r >{\raggedright\arraybackslash}X@{}}
\toprule
\textbf{Escenario} & \textbf{Vercel} & \textbf{Supabase} & \textbf{Mensual aprox.} & \textbf{Cuándo aplica} \\
\midrule
Piloto                  & Hobby (gratis) & Free  & \$0 MXN          & Validación inicial, 1 a 3 meses \\
Producción mínima       & Pro            & Free  & \textasciitilde\$360 MXN & Uso corporativo formal \\
Producción recomendada  & Pro            & Pro   & \textasciitilde\$810 MXN & Respaldos automáticos, soporte, SLA \\
\bottomrule
\end{tabularx}

\subsection*{Notas relevantes}

\begin{itemize}
  \item La operación es \textbf{técnicamente viable con ambos planes gratuitos}
        durante el piloto inicial.
        Las limitaciones principales son la ausencia de respaldos automáticos en Supabase Free
        y la cláusula de uso no comercial de Vercel Hobby.
  \item Si solo es posible adquirir un plan pago,
        \textbf{Vercel Pro tiene prioridad} sobre Supabase Pro:
        formaliza el uso comercial del hosting
        y habilita garantías sobre la capa pública de la aplicación.
        Supabase Pro puede sumarse después,
        cuando los respaldos automáticos o la capacidad lo justifiquen.
  \item Los costos cubren el stack completo,
        independientemente de cuántas herramientas se construyan sobre él.
\end{itemize}

A modo de contraste, una licencia de \emph{Zendesk Suite}
ronda los \$1{,}500 MXN mensuales por agente
y una licencia de \emph{Jira Standard} ronda los \$170 MXN mensuales por usuario.
La propuesta cubre a la organización completa
en el rango de \$0 a \$810 MXN mensuales totales.

% ============================================================
% 7. BENEFICIOS
% ============================================================
\section{Beneficios}

\begin{itemize}
  \item \textbf{Gobernanza centralizada.}
        Toda la infraestructura bajo cuentas corporativas
        administradas por el Área de Sistemas.
  \item \textbf{Trazabilidad operativa completa.}
        Cada acción queda auditada con autor, fecha y contexto;
        métricas operativas reales sobre las que hoy no se tiene visibilidad.
  \item \textbf{Reducción de fricción administrativa.}
        Alta y baja de usuarios sincronizadas automáticamente
        con Google Workspace; sin gestión manual de contraseñas ni accesos.
  \item \textbf{Aprovechamiento de capacidades internas.}
        Las herramientas se construyen, evolucionan y mantienen con personal interno,
        reduciendo la dependencia de proveedores externos para soluciones a la medida.
\end{itemize}

% ============================================================
% 8. RIESGOS Y MITIGACIONES
% ============================================================
\section{Riesgos y mitigaciones}

\begin{tabularx}{\textwidth}{@{}>{\raggedright\arraybackslash}p{5.5cm} >{\raggedright\arraybackslash}X@{}}
\toprule
\textbf{Riesgo} & \textbf{Mitigación} \\
\midrule
Dependencia de un solo desarrollador
  & Documentación técnica y operativa completa.
    Repositorio en organización corporativa.
    Stack popular, lo que facilita la sustitución o ampliación del equipo. \\
\addlinespace[0.3em]
Caída de Vercel o Supabase
  & Stack portable. Migración a infraestructura propia o a otro proveedor
    factible en plazos de semanas. \\
\addlinespace[0.3em]
\emph{Vendor lock-in}
  & Tecnologías estándar (PostgreSQL, Node.js).
    Sin componentes propietarios. Datos y código siempre exportables. \\
\addlinespace[0.3em]
Filtración o acceso indebido a datos
  & Reglas de acceso a nivel base de datos (RLS), HTTPS,
    sin contraseñas almacenadas, auditoría completa. \\
\addlinespace[0.3em]
Cuentas técnicas hoy a título personal
  & Plan de migración a cuentas corporativas previo
    a cualquier despliegue formal. Tema central de esta propuesta. \\
\addlinespace[0.3em]
Crecimiento de costos no anticipado
  & Escenarios de costos definidos. Escalado por etapas.
    Alertas configurables en planes pagos. \\
\bottomrule
\end{tabularx}

% ============================================================
% 9. PRÓXIMOS PASOS Y DECISIONES REQUERIDAS
% ============================================================
\section{Próximos pasos y decisiones requeridas}

Para avanzar, se requieren las siguientes definiciones por parte de la Gerencia de Sistemas:

\begin{enumerate}
  \item \textbf{Confirmación de dirección.}
        Aprobación para mover ambas herramientas
        a infraestructura corporativa bajo gobernanza de Sistemas.
  \item \textbf{Acceso al panel DNS (cPanel de HostGator).}
        Para configurar el subdominio
        \texttt{tickets.financieracrediflexi.com}
        mediante un registro CNAME, sin afectar el sitio web actual.
  \item \textbf{Coordinación con Google Workspace.}
        Creación de un proyecto en Google Cloud bajo el Workspace corporativo,
        para administrar el inicio de sesión único.
  \item \textbf{Correo corporativo administrativo.}
        Para registrar las cuentas de GitHub, Vercel y Supabase
        a nombre de la empresa.
  \item \textbf{Decisión sobre planes de pago.}
        Momento de transición de los planes gratuitos a planes Pro,
        con base en los escenarios de costos descritos.
  \item \textbf{Definición de gobernanza interna.}
        Personas del Área de Sistemas con acceso administrativo
        a las plataformas.
\end{enumerate}

\subsection*{Compromiso del solicitante}

\begin{itemize}
  \item Creación de todas las cuentas corporativas bajo los correos que se indiquen.
  \item Migración del repositorio y la base de datos a la infraestructura corporativa.
  \item Configuración del inicio de sesión con Google Workspace.
  \item Documentación técnica y operativa completa antes de la entrega formal.
  \item Soporte directo durante el periodo de validación inicial,
        con iteraciones basadas en el uso real de las herramientas.
\end{itemize}

% ============================================================
% GLOSARIO
% ============================================================
\section*{Glosario}
\addcontentsline{toc}{section}{Glosario}

\begin{description}[leftmargin=2.6cm, style=nextline, itemsep=0.3em]
  \item[CI/CD] Integración y despliegue continuos.
        Proceso automático que toma los cambios aprobados en el repositorio
        y los publica en producción.
  \item[CNAME] Tipo de registro DNS que apunta un subdominio
        (por ejemplo \texttt{tickets.financieracrediflexi.com}) hacia otro destino.
  \item[HTTPS] Protocolo de comunicación encriptada entre navegador y servidor.
  \item[JWT] Token criptográfico que prueba la identidad de un usuario autenticado
        en cada petición.
  \item[Next.js] Marco de desarrollo web sobre el que se construyen ambas aplicaciones.
        Mantenido por Vercel.
  \item[OAuth 2.0] Protocolo estándar de delegación de identidad.
        Es la base del inicio de sesión con Google.
  \item[PostgreSQL] Sistema gestor de base de datos relacional, abierto y estándar
        en la industria. Tecnología subyacente de Supabase.
  \item[RLS] \emph{Row Level Security}.
        Reglas de seguridad aplicadas a nivel base de datos,
        independientes del código de la aplicación.
  \item[SSO] \emph{Single Sign-On}, inicio de sesión único.
        Una sola credencial corporativa para acceder a todas las herramientas.
  \item[Vendor lock-in] Dependencia técnica de un proveedor específico
        que dificulta o encarece la migración hacia otro.
\end{description}

\end{document}
